% ============================================================
% 全国大学生数学建模竞赛 论文模板 (LaTeX)
% 2026 高教社杯 ｜ 参赛队伍专用
% 使用方式: 安装 TeX Live 或 MikTeX, 用 XeLaTeX 编译
%   中文支持: ctex 宏包 (自动处理中文字体)
% 编译命令:  xelatex main.tex  (至少编译两次以更新目录)
% ============================================================

\documentclass[12pt]{article}

% ---------- 中文字体与页面设置 ----------
\usepackage[UTF8]{ctex}
\usepackage[a4paper, top=2.5cm, bottom=2.5cm, left=2.5cm, right=2.5cm]{geometry}
\usepackage{amsmath, amssymb, amsfonts}
\usepackage{graphicx}
\usepackage{booktabs}      % 三线表
\usepackage{multirow}
\usepackage{array}
\usepackage{enumitem}
\usepackage{algorithm}
\usepackage{algorithmic}
\usepackage{caption}
\usepackage{subcaption}
\usepackage{xcolor}  % xcolor 而非 color：lstlisting 的 gray 颜色仅 xcolor 预定义
\usepackage{float}
\usepackage{hyperref}
\usepackage{setspace}
\usepackage{fancyhdr}
\usepackage{listings}      % 代码排版

% ---------- 页眉页脚 ----------
\pagestyle{fancy}
\fancyhf{}
\fancyhead[L]{\small 队伍编号：XXXXX}
\fancyhead[R]{\small 全国大学生数学建模竞赛}
\fancyfoot[C]{\thepage}
\renewcommand{\headrulewidth}{0.4pt}

% ---------- 超链接样式 ----------
\hypersetup{
  colorlinks=true,
  linkcolor=black,
  citecolor=black,
  urlcolor=blue
}

% ---------- 代码样式 ----------
\lstset{
  basicstyle=\ttfamily\small,
  numbers=left,
  numberstyle=\tiny\color{gray},
  frame=single,
  breaklines=true,
  language=Python
}

% ---------- 标题设置 ----------
\title{\textbf{题 目：请在这里填写论文题目}}
\author{队伍编号：XXXXX}
\date{}

\begin{document}

% ============================================================
% ⚠️ 官方 2026 修订稿规定（见 references/04_论文写作/05_论文格式规范.md）：
%   - 第 1 页承诺书、第 2 页编号专用页均由组委会提供，**不要自制封面页**，
%     且论文任何地方不得出现姓名/学校/赛区信息（第六条）；
%   - 电子版论文第一页必须为摘要专用页（第十条）；
%   - 正文**不要目录**（第四条）。
% 因此本模板从摘要专用页开始，页码从 1 起。
% ============================================================

% ============================================================
% 摘要专用页（第 1 页；重要！占评分 20%-25%，最后写）
% ============================================================
\begin{center}
{\LARGE\bfseries 论文题目（居中，一句话概括核心工作）\par}
\vspace{0.6cm}
{\LARGE\bfseries 摘 \quad 要}
\end{center}

% —— 摘要正文从这里开始（约一页，含每问方法 + 数字结果）——
本文针对\underline{\hspace{4cm}}问题，建立了\underline{\hspace{3cm}}模型。

针对问题一，本文采用\underline{\hspace{3cm}}方法，建立\underline{\hspace{3cm}}模型，
得到\underline{\hspace{3cm}}结果（关键数字：\underline{\hspace{3cm}}）。

针对问题二，本文构建了\underline{\hspace{3cm}}模型，求解得\underline{\hspace{3cm}}
（具体数值结果，如最优时点、最优方案）。

针对问题三，在问题二基础上引入\underline{\hspace{3cm}}因素，
建立\underline{\hspace{3cm}}模型，得到\underline{\hspace{3cm}}。

针对问题四，本文提出\underline{\hspace{3cm}}判定方法，
阈值设为\underline{\hspace{3cm}}，灵敏度\underline{\hspace{3cm}}，
特异度\underline{\hspace{3cm}}。

最后，对模型进行了灵敏度分析（关键参数 $\pm 5\%$、$\pm 10\%$ 变化，
结果波动\underline{\hspace{2cm}}），表明模型具有较好的鲁棒性。

\vspace{0.5cm}
\noindent\textbf{关键词：}\underline{\hspace{2cm}}；\underline{\hspace{2cm}}；\underline{\hspace{2cm}}；\underline{\hspace{2cm}}；\underline{\hspace{2cm}}

\newpage
% ⚠️ 官方规定正文"不要目录"（第四条），此处不设 \tableofcontents。
% 论文题目写在摘要专用页顶部：

% ============================================================
% 一、问题重述与分析
% ============================================================
\section{问题重述与分析}
\subsection{问题背景}
% 用自己的话描述问题背景（不要复制粘贴题目原文）

\subsection{问题重述}
% 每个小问用一句话概括
\begin{enumerate}[label=问题\arabic*：]
  \item \underline{\hspace{6cm}}
  \item \underline{\hspace{6cm}}
  \item \underline{\hspace{6cm}}
  \item \underline{\hspace{6cm}}
\end{enumerate}

\subsection{问题分析}
% 给出整体解题思路，建议配流程图
% \begin{figure}[H]
%   \centering
%   \includegraphics[width=0.8\textwidth]{flowchart.png}
%   \caption{解题流程图}
% \end{figure}

% ============================================================
% 二、模型假设与符号说明
% ============================================================
\section{模型假设与符号说明}
\subsection{模型假设}
% 3-5 条，每条 = 假设 + 理由
\begin{enumerate}[label=假设\arabic*：]
  \item \underline{\hspace{5cm}}。理由：\underline{\hspace{4cm}}。
  \item \underline{\hspace{5cm}}。理由：\underline{\hspace{4cm}}。
  \item \underline{\hspace{5cm}}。理由：\underline{\hspace{4cm}}。
\end{enumerate}

\subsection{符号说明}
% 三线表：符号 + 含义 + 单位
\begin{table}[H]
\centering
\caption{主要符号说明}
\begin{tabular}{ccc}
\toprule
符号 & 含义 & 单位 \\
\midrule
$x_i$ & \underline{\hspace{3cm}} & \underline{\hspace{1cm}} \\
$y_j$ & \underline{\hspace{3cm}} & \underline{\hspace{1cm}} \\
$\alpha$ & \underline{\hspace{3cm}} & \underline{\hspace{1cm}} \\
\bottomrule
\end{tabular}
\end{table}

% ============================================================
% 三、数据预处理（如有数据）
% ============================================================
\section{数据预处理}
\subsection{数据来源与概况}
% 样本量、变量数、数据类型

\subsection{缺失值与异常值处理}
% 处理方法 + 理由

\subsection{特征工程}
% 哑变量、衍生变量、标准化

% ============================================================
% 四、问题一的模型建立与求解
% ============================================================
\section{问题一的模型建立与求解}
\subsection{模型建立}
% 思路 + 公式。公式用 equation 环境，全文自动连续编号，不按章节编号；
% 公式前一句文字引出，公式后逐符号解释并给参数取值，后文用 式~\eqref{eq:model1} 回引。
\begin{equation}
y = \beta_0 + \beta_1 x + \varepsilon
\label{eq:model1}
\end{equation}
式中：$\beta_0$ 为【含义，取值___】；$\beta_1$ 为【含义，取值___】。
如式~\eqref{eq:model1} 所示，……（正文引用公式一律用编号回引）。

\subsection{模型求解}
% 求解方法、工具与版本（如 YALMIP+Gurobi, MIPGap=1e-4）、关键代码（放附录）

\subsection{结果与分析}
% 结果表格 + 图 + 2-3 句解读；图注可携带情景参数
% 双图并排示例（需 graphicx + subcaption，已加载）：
% \begin{figure}[H]
%   \centering
%   \begin{subfigure}[b]{0.48\textwidth}
%     \centering
%     \includegraphics[width=\textwidth]{fig_a.png}
%     \caption{最小螺距-次数图}
%   \end{subfigure}
%   \hfill
%   \begin{subfigure}[b]{0.48\textwidth}
%     \centering
%     \includegraphics[width=\textwidth]{fig_b.png}
%     \caption{最小螺距-区间图}
%   \end{subfigure}
%   \caption{改进二分查找收敛过程}
% \end{figure}

% ============================================================
% 五、问题二的模型建立与求解
% ============================================================
\section{问题二的模型建立与求解}
\subsection{模型建立}
% 优化三段式：决策变量 → 目标函数 → 逐条约束（每条后跟一句物理意义），最后汇总为 s.t. 块
\begin{align}
\min \quad & f(\mathbf{x}) \\
\text{s.t.} \quad & g_i(\mathbf{x}) \le 0, \quad i = 1, \dots, m, \nonumber \\
& x_j \ge 0, \quad j = 1, \dots, n \nonumber
\end{align}

\subsection{模型求解}
% 算法流程（伪代码）
\begin{algorithm}[H]
\caption{求解算法}
\begin{algorithmic}[1]
\STATE 初始化参数
\WHILE{未收敛}
  \STATE 更新变量
\ENDWHILE
\RETURN 最优解
\end{algorithmic}
\end{algorithm}

\subsection{结果与分析}

% ============================================================
% 六、问题三的模型建立与求解
% ============================================================
\section{问题三的模型建立与求解}
\subsection{模型建立}
\subsection{模型求解}
\subsection{结果与分析}

% ============================================================
% 七、问题四的模型建立与求解
% ============================================================
\section{问题四的模型建立与求解}
\subsection{模型建立}
\subsection{模型求解}
\subsection{结果与分析}

% ============================================================
% 八、模型评价与改进
% ============================================================
\section{模型的评价、改进与推广}
\subsection{灵敏度分析}
% 关键参数 ±5%、±10% 变化对结果的影响
% 用表格呈现：参数变化 -> 结果变化 -> 是否稳定

\subsection{模型的优点}
% 具体、客观

\subsection{模型的缺点与改进方向}
% 说得"无关紧要"些：数据误差、假设简化

\subsection{模型的推广}
% 一句话：可用于哪些类似场景

% ============================================================
% AI 工具使用声明（2026 新规硬性要求：必须在【参考文献之前】！）
% ============================================================
\section*{AI 工具使用声明}
本参赛队在竞赛过程中使用了 AI 工具，主要用于【语言润色、代码调试、资料查询】等辅助工作，详细使用情况见支撑材料。
% 若未使用 AI 工具，将上一行整段替换为：
%   本参赛队在竞赛过程中未使用任何 AI 工具。
% 使用 AI 的队还需另附支撑材料 PDF：AI工具使用详情.pdf
%   内容含：工具名称与版本、使用环节、提示方式、人工核验情况

% ============================================================
% 参考文献 (GB/T 7714)
% ============================================================
\begin{thebibliography}{99}
\bibitem{ref1} 作者. 题目[J]. 期刊名, 年, 卷(期): 页码.
\bibitem{ref2} 作者. 书名[M]. 出版地: 出版社, 年份.
\bibitem{ref3} 作者. 文章题目[C]//会议论文集. 出版地: 出版社, 年份: 页码.
\end{thebibliography}

% ============================================================
% 附录：核心代码
% ============================================================
\appendix
\section{核心代码}
\begin{lstlisting}
# 在这里贴核心代码（不放正文）
import numpy as np
import pandas as pd
\end{lstlisting}

\end{document}
